
défini par \(f_n(X_i)=X_i\) si \(i<n\), \(f_n(X_n)=0\). Il se factorise par l'inclusion de \(K[X_1,\ldots,X_{n-1}]\) dans \(K[X_1,\ldots,X_n]\), ce qui montre que l'homomorphisme de \(H(K[X_1,\ldots,X_{n-1}])\) dans \(H(K[X_1,\ldots,X_n])\) est injectif, et nous permettra d'identifier le permier à un sous-groupe du second, et de même nous considérerons \(H(K[X_1,\ldots,X_n])\) comme contenu dans \(H(K[X_i]_{i\geq 1})\).

On notera que \(f_1\) n'est autre que l'augmentation de \(K[X]\).

1.5. On plonge \(K[X_1,\ldots,X_n]\) dans \(K[T_1,\ldots,T_n]\) en envoyant \(X_i\) sur la \(i\)-ème fonction symétrique des \(T_j\)~; en faisant opérer le groupe symétrique \(\mathfrak S_n\) par permutation des \(T_j\), et en utilisant le théorème des polynômes symétriques et (1.3.ii)), on voit que l'homomorphisme de \(H(K[X_1,\ldots,X_n])\) dans \(H(K[T_1,\ldots,T_n])\) qu'on en déduit est injectif, ce qui nous permettra encore d'identifier le premier à un sous-groupe du second~:
\[
  H(K[X_1,\ldots,X_n])\simeq H(K[T_1,\ldots,T_n]^{\mathfrak S_n})\xrightarrow{\sim} H(K[T_1,\ldots,T_n])^{\mathfrak S_n}\subset H(K[T_1,\ldots,T_n])
\]% source-fix p300 §1.5 (2026-07-03): restored Grothendieck's text vs the paraphrase — removed interpolated «élémentaire», «.»→«;», «qui s'en déduit»→«qu'on en déduit», «permet»→«nous permettra», inlined the H(...) terms, restored dropped intermediate iso H(K[T]^{Sym_n})→~H(K[T])^{Sym_n}.
On note
\[
  u_i^n:K[X]\longrightarrow K[T_1,\ldots,T_n]
\]
l'homomorphisme défini par \(u_i^n(X)=T_i\), et on pose
\[
  u_H^n=\sum_{i=1}^n H(u_i^n):H(K[X])\longrightarrow H(K[T_1,\ldots,T_n]).
\]
Le groupe symétrique permute les morphismes \(H(u_i^n)\), et laisse donc invariant le morphisme \(u_H^n\), qui par suite prend ses valeurs dans
